\documentclass[10pt,a4paper]{amsart}
\usepackage{amsmath,amssymb,amsthm,graphicx}
% Minimal class-compatibility preamble for a reader excerpt of
%   Juan Esteban Rodriguez Camargo, "An integral model of the perfectoid modular curve",
%   Journal de l'Ecole polytechnique -- Mathematiques 8 (2021), DOI 10.5802/jep.170.
%
% The block below is the author's own macro and theorem declaration block of the pinned
% source file, copied mechanically and unchanged (source lines 6-91), followed by the
% abbreviations that the publisher class cedram supplies and that the retained author text
% uses (same meanings as in that class). No mathematical macro, statement, symbol or
% assumption is introduced, changed or removed by this file.
\usepackage{xspace}

\usepackage[scr=rsfs,cal=euler]{mathalfa}
\multlinegap0pt

\newenvironment{enumeratei}
{\bgroup\def\theenumi{\roman{enumi}}\def\theenumii{\arabic{enumii}}\begin{enumerate}}
{\end{enumerate}\egroup}

\newcommand\bw[1]{\mathbin{\hbox{$\mathsf{\Lambda}$}^{\!#1}}}
\newcommand{\quand}{\quad\text{and}\quad}
\newcommand{\sbullet}{{\scriptscriptstyle\bullet}}

\newcommand\mto{\mathchoice{\longmapsto}{\mapsto}{\mapsto}{\mapsto}}
\newcommand\hto{\mathrel{\lhook\joinrel\to}}
\newcommand\To[1]{\mathchoice{\xrightarrow{\textstyle\kern4pt#1\kern3pt}}{\stackrel{#1}{\to}}{}{}}
\let\ra\rightarrow

\let\ts\textstyle
\newcommand{\lcr}{[\![}
\newcommand{\rcr}{]\!]}
\newcommand{\lpr}{(\!(}
\newcommand{\rpr}{)\!)}

\newcommand{\rae}{\mathrm{ae}}
\newcommand{\an}{\mathrm{an}}
\newcommand{\rb}{\mathrm{b}}
\newcommand{\rc}{\mathrm{c}}
\newcommand{\qc}{\mathrm{qc}}
\newcommand{\cusp}{\mathrm{cusp}}
\newcommand{\cusps}{\mathrm{cusps}}
\newcommand{\cyc}{\mathrm{cyc}}
\newcommand{\et}{\mathrm{\acute{e}t}}
\newcommand{\fTd}{\mathrm{fTd}}
\newcommand{\pr}{\mathrm{pr}}
\newcommand{\proet}{\mathrm{pro}\textup{-}\et}
\newcommand{\sm}{\mathrm{sm}}
\newcommand{\univ}{\mathrm{univ}}

\newcommand{\wotimes}{\mathbin{\widehat\otimes}}

\usepackage[all,cmtip]{xy}\let\labelstyle\textstyle
\usepackage{tikz-cd}

\numberwithin{equation}{section}

\newtheorem{prop}{Proposition}[section]
\newtheorem{lem}[prop]{Lemma}
\newtheorem{theo}[prop]{Theorem}
\newtheorem{cor}[prop]{Corollary}

\theoremstyle{definition}
\newtheorem*{claim}{Claim}
\newtheorem{exam}[prop]{Example}
\newtheorem{definition}[prop]{Definition}
\newtheorem{remark}[prop]{Remark}

\newcommand{\bb}[1]{\mathbb{#1}}
\newcommand{\n}[1]{\mathcal{#1}}
\newcommand{\f}[1]{\mathfrak{#1}}
\newcommand{\s}[1]{\mathscr{#1}}

\DeclareMathOperator{\GL}{GL}
\DeclareMathOperator{\Tr}{Tr}
\DeclareMathOperator{\id}{id}
\DeclareMathOperator{\SL}{SL}
\DeclareMathOperator{\R}{R}
\DeclareMathOperator{\D}{D}
\DeclareMathOperator{\N}{N}

\DeclareMathOperator{\Hom}{Hom}
\DeclareMathOperator{\dotimes}{\otimes^{\mathrm{L}}}
\DeclareMathOperator{\coker}{coker}

\DeclareMathOperator{\Gal}{Gal}
\DeclareMathOperator{\Spf}{Spf}
\DeclareMathOperator{\Spec}{Spec}
\DeclareMathOperator{\Ell}{Ell}
\DeclareMathOperator{\Spa}{Spa}
\DeclareMathOperator{\Cusps}{Cusps}
\DeclareMathOperator{\Tate}{Tate}
\DeclareMathOperator{\sHom}{\s{H}\!\mathit{om}}
\DeclareMathOperator{\HomSurj}{HomSurj}
\DeclareMathOperator{\Art}{Art}
\DeclareMathOperator{\Intperf}{\f{I}nt\textup{-}\f{P}erf}
\DeclareMathOperator{\Perf}{\s{P}erf}
\newcommand{\mc}[1]{X(Np^{#1})}
\newcommand{\fmc}[1]{\mathfrak{X}(Np^{#1})}

% Publisher-class abbreviations (cedram) used by the retained author text.
\providecommand{\loccit}{\emph{loc. cit.}}
\providecommand{\bibliofont}{\small}
\providecommand{\backmatter}{}
\providecommand{\skpt}{\smallskip}
\providecommand{\smaller}{\small}
\newcommand{\ie}{i.e.,\xspace}
\newcommand{\eg}{e.g.,\xspace}
\newcommand{\etc}{etc.\xspace}
\newcommand{\cf}{cf.\xspace}
\newcommand{\resp}{resp.\xspace}
\title{Local Serre-Pontryagin duality for a locally free sheaf on an lci projective curve over the valuation ring of a finite extension of Q\_p}
\author{Juan Esteban Rodríguez Camargo}
\date{doi:10.5802/jep.170; CC BY 4.0}
\begin{document}
\maketitle
\setcounter{section}{3}
\setcounter{prop}{4}
\noindent\textit{Editorial note.} The counted unit is the author's Proposition 3.7 at source lines 732-745, source label LocalDualityCurves, together with the proof environment that begins two lines later at 747-763; the intervening line 746 is blank in the source and is not retained. The printed number is fixed by initial\_counters so that the counted proposition carries its original number 3.7: the retained excerpt also contains the author's Definition at 649 (original number 3.1) and the quoted Hartshorne theorem at 659 (original number 3.2), which the excerpt prints as 3.5 and 3.6 because the items that separate them in the published article are not retained.\par\medskip
\noindent\textit{Editorial note.} The proof span at source lines 659-673 is the author's own quotation of Grothendieck-Serre duality (Theorem 3.2, label GrothSerreDuality) together with the author's citation proof pointer; it is retained because the counted proof at source line 758 invokes that theorem by its label, and it is not counted as a second proposition.\par\medskip
\noindent\textit{Editorial note.} The context span at source lines 645-657 is the section opening with the author's derived-category notation and the definition of embeddable and lci morphisms; the context span at 730 is the author's own sentence introducing the local duality theorem. Neither retained passage was altered.\par\medskip
\noindent\textit{Editorial note.} The retained preamble reproduces the author's macro and theorem block verbatim from source lines 6-91 and adds, as a comment-documented adapter, the abbreviations that the publisher class cedram supplies and that the retained sentences use (i.e., e.g., cf., etc., resp., loc. cit., smallskip and backmatter); no mathematical macro, symbol, statement or hypothesis is added, changed or removed. The cedram publisher class is not installed on the assembly host, so the wrapper class amsart is used.\par\medskip
\noindent\textit{Editorial note.} The retained bibliography carries exactly the two works cited by the retained text (Hartshorne, Residues and duality; Mumford, Abelian varieties), copied verbatim from the public corpus record of this same source version.\par\medskip
\noindent\textit{Editorial scope.} One original result: Proposition 3.7 (label LocalDualityCurves) of the source with its complete author proof at source lines 747-763 and the author's own retained quotation of Grothendieck-Serre duality that this proof invokes. The already-counted Theorem 2.8 of the same source and the later vanishing, duality and completed-cohomology theorems are not counted here and are not used as premises.\par\medskip
\noindent\textit{Allowed standard prerequisites (editorial).}\begin{itemize}
\item Grothendieck-Serre duality for projectively embeddable morphisms. For a projectively embeddable morphism f from X to Y of noetherian schemes of finite Krull dimension there exist an exceptional inverse image functor f\textasciicircum{}!, a trace map Tr from Rf\_* f\textasciicircum{}! to 1 in D\textasciicircum{}+\_qc(Y) and the adjunction theta from Rf\_* R-Hom\_X(F, f\textasciicircum{}!G) to R-Hom\_Y(Rf\_*F, G); the author restates exactly this as Theorem 3.2 at source lines 659-669 and its retained proof at 671-673 is the author's own pointer to Hartshorne, Residues and duality, Th. III.8.7, III.10.5 and III.11.1.. Supplies the quasi-isomorphism Rf\_*(F-dual tensored with omega\textasciicircum{}circ)[1] = R-Hom over O\_K of (Rf\_*F, O\_K) used at source lines 758-761; the counted proof does not reprove it.
\item Pontryagin duality over the valuation ring. For the discrete valuation ring O\_K with fraction field K and a finite O\_K-module M, the natural map from M to Hom over O\_K of (Hom over O\_K of (M, K/O\_K), K/O\_K) is an isomorphism; in the notation of the counted proof this reads Hom over O\_K of (ker(d tensored with 1), K/O\_K) = coker(d-dual).. Turns the identification of the two displayed two-term complexes into the perfect pairing asserted by the counted proposition (source line 762).
\end{itemize}\medskip
Let $K$ be a finite extension of $\bb{Q}_p$ and $\n{O}_K$ its valuation ring. In this section we recall the Grothendieck-Serre duality theorem for local complete intersection (lci) projective curves over $\n{O}_K$, we will follow \cite{hartshorne1966residues}. Then, we use Pontryagin duality to define a local duality paring of coherent cohomologies.

Let $X$ be a locally noetherian scheme and $\D(X)$ the derived category of $\s{O}_X$\nobreakdash-mod\-ules. We use subscripts $\rc, \qc$ on $\D(X)$ for the derived category of $\s{O}_{X}$\nobreakdash-modules with coherent and quasi-coherent cohomology, the subscript $\fTd$ refers to the subcategory of complexes with finite Tor dimension. We use superscripts $+,-,\rb$ for the derived category of bounded below, bounded above and bounded complexes respectively. For instance, $\D_{\rc}^b(X)_{\fTd}$ is the derived category of bounded complexes of $\s{O}_{X}$-modules of finite Tor dimension and coherent cohomology. If $X=\Spec A$ is affine, we set $\D(A):=\D_{\qc}(X)$, the derived category of $A$-modules.

\begin{definition}
Let $f:X\to Y$ be a morphism of schemes.
\begin{enumerate}
\item The map $f$ is \textit{embeddable} if it factors as $X\To{\iota} S \to Y $ where $\iota$ is a finite morphism and $S$ is smooth over $Y$.
\item The map $f$ is \textit{projectively embeddable} if it factors as composition \hbox{$ X\To{\iota} \bb{P}^n_Y\to Y$} for some $n\geq 0$, where $\iota$ a finite morphism.

\item The map $f$ is a \textit{local complete intersection} if locally on $Y$ and $X$ it factors as $X\To{\iota} S \to Y $, where $S$ is a smooth $Y$-scheme, and $\iota$ is a closed immersion defined by a regular sequence of $S$. The length of the regular sequence is called the codimension of $X$ in $S$.
\end{enumerate}
\end{definition}

\begin{theo}[Hartshorne]
\label{GrothSerreDuality}

Let $f: X\to Y$ be a projectively embeddable morphism of noetherian schemes of finite Krull dimension. Then there exist an exceptional inverse image functor $f^!: \mathrm{D}(Y)\to \mathrm{D}(X)$, a trace map $\Tr: \mathrm{R}f_*f^!\to 1$ in $\mathrm{D}^+_{\qc}(Y)$, and an adjunction
\[
\theta: \R f_* \R \sHom_{X}(\n{F}, f^! \n{G}) \To{\sim} \R\sHom_Y(\R f_* \n{F}, \n{G})
\]
for $\s{F}\in \mathrm{D}^-_{\qc}(X)$ and $\s{G}\in \mathrm{D}^+_{\qc}(Y)$.

Moreover, the formation of the exceptional inverse image is functorial. More precisely, given a composition $X\To{f} Y \To{g} Z$ with $f,g$ and $gf$ projectively embeddable, there is a natural isomorphism $(gf)^!\cong f^!g^!$. This functor commutes with flat base change. Namely, let $u:Y'\to Y$ be a flat morphism, $f':X'\to Y'$ the base change of $X$ to $Y'$ and $v: X'\to X$ the projection. Then there is a natural isomorphism of functors $v^*f^!={f'}^! u^*$.
\end{theo}

\begin{proof}
We refer to \cite[Th.\,III.8.7]{hartshorne1966residues} for the existence of $f^!$, its functoriality and compatibility with flat base change. See Theorems III.10.5 and III.11.1 of \loccit\ for the existence of $\Tr$ and the adjunction $\theta$ respectively.
\end{proof}

We now prove the local duality theorem for vector bundles over lci projective curves:

\begin{prop}
\label{LocalDualityCurves}
Let $f: X\to \Spec \n{O}_K$ be an lci projective curve, and let $\omega_{X/\n{O}_K}^\circ$ be the dualizing sheaf of $f$, \ie the invertible sheaf such that $\omega_{X/\n{O}_K}^\circ[1]=f^!\n{O}_K$. Let~$\n{F}$ be a locally free $\s{O}_X$-module of finite rank, then:
\begin{enumerate}

\item $\R f_*\n{F}$ is representable by a perfect complex of length $[0,1]$;

\item we have a perfect pairing
\begin{align*}
\mathrm{H}^0(X,\n{F}\otimes K/\n{O}_K)\times \mathrm{H}^1(X, \n{F}^\vee\otimes \omega_{X/\n{O}_K}^\circ)\to K/\n{O}_K
\end{align*}
\end{enumerate}
given by the composition of the cup product and the trace $\Tr: Rf_* \omega^{\circ}_{X/\n{O}_K}\to \n{O}_K$.
\end{prop}

\begin{proof}
As $\n{F}$ is a vector bundle and $f$ is projective of relative dimension $1$, the cohomology groups $\R^i f_* \n{F}$ are finitely generated over $\n{O}_K$ and concentrated in degrees~$0$ and~$1$. Then, $\R f_* \n{F}$ is quasi-isomorphic to a complex
\[
0 \to M_0 \To{d}M_1\to 0,
\]
with $M_1$ and $M_2$ finite free $\n{O}_K$-modules. Moreover, the complex
\[
0 \to M_0\otimes K/\n{O}_K \To{d\otimes 1} M_1\otimes K/\n{O}_K \to 0
\]
is quasi-isomorphic to $\R f_*(\n{F}\otimes K/\n{O}_K)$ in $\D(\n{O}_K)$, see \cite[Th.\,5.2]{mumford1974abelian}.

The duality theorem \ref{GrothSerreDuality} gives a quasi-isomorphism
\[
\R f_*(\n{F}^{\vee}\otimes \omega^{\circ}_{X/\n{O}_K})[1] =\R f_*\R\sHom_{X}(\n{F}, f^!\n{O}_K)\simeq \R\Hom_{\n{O}_K}(\R f_*\n{F}, \n{O}_K).
\]
This implies that $\R f_* (\n{F}^\vee \otimes \omega_{X/\n{O}_K}^\circ)$ is quasi-isomorphic to $0 \to M_1^\vee \To{d^\vee}M_0^\vee \to 0$. Finally, Pontryagin duality for $\n{O}_K$ implies $\Hom_{\n{O}_K}(\ker(d\otimes 1), K/\n{O}_K)= \coker d^\vee$, which translates in the desired statement.
\end{proof}

\begin{thebibliography}{99}
\bibitem{hartshorne1966residues} R. Hartshorne - Residues and duality , Lect. Notes in Math. , vol. 20 , Springer-Verlag, Berlin-New York, 1966
\bibitem{mumford1974abelian} D. Mumford - Abelian varieties , TIFR Studies in Mathematics , vol. 5 , Tata Institute of Fundamental Research, Bombay, 2008
\end{thebibliography}
\end{document}
